\section{Related Work}
\label{sec:related}

\paragraph{Benchmarks for computer-use agents.} Interactive benchmarks evaluate CUAs on desktop operating systems \citep{xie2024osworld,bonatti2025windowsagentarena,yang2025macosworld,davydova2025osuniverse}, websites \citep{zhou2024webarena,koh2024visualwebarena,he2024webvoyager,xue2025illusion}, enterprise software \citep{drouin2024workarena,boisvert2024workarenapp,xu2024theagentcompany} and mobile apps \citep{rawles2025androidworld}, alongside offline trajectory datasets \citep{deng2023mind2web} and grounding tests \citep{li2025screenspotpro}. They judge the final state with deterministic checkers or rubric-based LLM judges \citep{xue2025illusion,lu2025agentrewardbench}, and they differ mainly in domain and horizon, up to tasks that take a person more than an hour \citep{yuan2026osworld2}. Time enters only as a step cap, and the cap itself shifts scores \citep{xlang2025osworldverified,qin2025uitars,agashe2025agents2}; AndroidWorld derives its per-task cap from human step counts \citep{rawles2025androidworld}. Robustness benchmarks add pop-ups and other everyday disruptions \citep{zhang2025popups,yang2026guirobust}, and PersonaCUABench finds that many fully completed runs still fall short of the requesting persona's satisfaction \citep{li2026personacuabench}. \bench keeps deterministic end-state checks, but a correct end state is not enough: the run must also stay within a human-labelled tolerance.

\paragraph{Efficiency, speed and reliability.} A growing body of work argues that accuracy alone misrepresents agents. Accounting for cost \citep{kapoor2025agentsthatmatter,kapoor2026hal}, consistency across repeated runs \citep{yao2025taubench,rabanser2026reliability} and benchmark rigour \citep{zhu2025abc} each changes the conclusions drawn from pass rates. For CUAs, OSWorld-Human annotates a human trajectory for every OSWorld task and finds that the best agents take 2.7--4.3$\times$ more steps than necessary, with end-to-end latencies of tens of minutes \citep{abhyankar2026osworldhuman}. cua-speedrun standardizes the infrastructure for timing CUAs and finds that frontier agents spend 86--127\,s per OSWorld task \citep{aggarwal2026cuaspeedrun}. ComponentBench pairs 2,910 programmatically generated UI-component tasks with human reference trajectories, and even its fastest agent configuration takes 3.7$\times$ as long as the matched human \citep{guan2026componentbench}. METR measures agent capability by the human completion time of the tasks agents can do \citep{kwa2025metr}, and ARC-AGI-3 scores agents by their action efficiency relative to humans on abstract games \citep{arcprize2026arcagi3}. These works use human effort as a reference for efficiency or difficulty. We also ask annotators how much more than their own effort they would tolerate from an agent, and we use that tolerance as a pass/fail boundary on realistic, human-written tasks.

\paragraph{Easy for people, hard for agents.} Benchmarks made of tasks that people solve easily expose gaps that hard tasks hide, including web games \citep{thomas2025webgames}, web browsing questions \citep{song2025bearcubs}, abstract reasoning puzzles \citep{chollet2019measure,chollet2025arcagi2,arcprize2026arcagi3} and commonsense questions \citep{simplebench2024}. The pattern echoes Moravec's paradox \citep{moravec1988mind}: the perception and motor skills that people exercise without effort are among the hardest to automate. \bench applies this design to ordinary interactions in ordinary applications, and asks both whether agents succeed at them and whether they do so before a person would give up.

\paragraph{Failure modes of agents.} Taxonomies of agent failures \citep{cemri2025mast,zhu2025agentdebug} and studies of overthinking, in which reasoning models spend far more computation on easy problems than they need \citep{chen2025overthinking,cuadron2025danger}, explain why agents fail or waste effort. For GUI agents, repetitive loops and premature stopping recur as dominant failure modes \citep{koh2024visualwebarena,han2026vlaagui,wei2026steplevel}, and AgentRewardBench labels repetitiveness alongside success and side effects \citep{lu2025agentrewardbench}. These taxonomies are organized around the agent's internal errors. Our frustration patterns (\S\ref{sec:patterns}) are organized around what the user sees, and they include patterns such as stalling and detours that a correctness criterion does not penalize at all.

\paragraph{User frustration and tolerance.} HCI research defines frustration as the appraisal of an interface as an obstacle to a goal \citep{bessiere2006model}, and diary studies found that it costs users a third to a half of their time at the computer \citep{ceaparu2004frustration,lazar2006severity}. Waiting is a central cause. Classic response-time limits place the end of uninterrupted attention at about ten seconds \citep{miller1968response,card1991information,nielsen1993usability}, web users tolerate about two seconds when retrieving information \citep{nah2004tolerable}, and task performance and the intention to return drop with delays of a few seconds \citep{galletta2004delays}. Expectations come from how long a task usually takes, so a result that takes far longer than expected frustrates \citep{shneiderman1984response}, and direct-manipulation tasks require faster responses than conversational ones \citep{dabrowski2011forty}. Human effort can be predicted from operator times with the Keystroke-Level Model and GOMS \citep{card1980keystroke,card1983psychology}; we measure it directly instead. People stop relying on automation after they observe it err \citep{parasuraman1997humans,lee2004trust,dietvorst2015algorithm}, and guidelines for human-AI interaction call for efficient dismissal and correction of an assistant \citep{horvitz1999principles,amershi2019guidelines}. Recent studies of LLM agents find that latency is a main complaint and a reason users take over \citep{mozannar2025magentic,maslych2025mitigating}, that deployed agents loop and run slower than doing the task by hand \citep{shome2026johnny,zhang2025unintended}, and that oversight depends on when control passes between user and agent \citep{huq2025cowpilot,huq2026modeling,chen2026oversight}. Tolerance for delay also depends on the task \citep{tan2026latency}. We therefore label tolerance per task instead of assuming a global limit, and we use it as the boundary between an acceptable and a frustrating run.
